\section{Mechanism}
\label{sec:mechanism}

Aggregate accuracy does not explain itself. This section traces one task end to
end, then reports where the treatment arm fails.

\subsection{Possession is not delivery}

Task 1278 asks for \emph{the average fee NexPay would charge for a credit
transaction of 50 EUR}.

\begin{center}
\footnotesize
\setlength{\tabcolsep}{4pt}
\begin{tabular}{@{}lll@{}}
\toprule
 & \tabhead{Answer} & \\
\midrule
Gold      & \texttt{0.352294} & \\
\armC{}   & \texttt{0.352294} & correct \\
\armM{}   & \texttt{0.353053} & wrong \\
\armS{}   & \texttt{0.353053} & wrong \\
\bottomrule
\end{tabular}
\end{center}

Both figures reproduce directly against the warehouse:

\begin{center}
\footnotesize
\setlength{\tabcolsep}{4pt}
\begin{tabular}{@{}ll@{}}
\toprule
\texttt{WHERE is\_credit IS NULL OR is\_credit = true} & \texttt{0.352294} \\
\texttt{WHERE is\_credit = true} & \texttt{0.353053} \\
\bottomrule
\end{tabular}
\end{center}

A \texttt{NULL} in a \texttt{fees} column means \emph{this rule applies to
every value}, not \emph{unknown}. Both baselines dropped the \texttt{NULL} rows
and produced the identical wrong number. The contract arm called
\texttt{lookup\_domain("fees")}, and its recorded reasoning states the rule
back: \emph{``rules that apply to credit transactions: is\_credit true or
null''}.

\textbf{The decisive detail is that \texttt{manual.md} documents this rule
too.} \armM{} had it --- inside a 24{,}177-character system prompt. \armC{}'s
resident prompt is 2{,}475 characters, with the rule reachable by one targeted
call.

The finding is not that the contract arm was given more information. Both arms
had the information. Only one of them used it. What the contract changes is not
what is knowable but what is \emph{reachable at the moment of use} --- which is
a claim about delivery, and it is why a result that looks like it is about
documentation is really about retrieval structure over the same documentation.

\subsection{Where the treatment arm loses}

In run~A, \armC{} loses 14 tasks to \armS{} and 29 to \armM{}, 10 to both. Two
modes account for nearly all of them.

\textbf{Set membership under \texttt{NULL} wildcards} --- the same rule as
above, failing in the other direction. On task 1500 (\emph{fee IDs that apply
to account\_type = O and aci = C}) \armC{} returns both false positives and
false negatives while \armM{} matches exactly. Applying ``\texttt{NULL} means
all'' to an enumeration is harder than applying it to an average: the model
must decide, per column, whether a \texttt{NULL} widens the match.

\textbf{Near-miss arithmetic} --- task 1274, gold \texttt{0.126459}, \armC{}
\texttt{0.125516}: a real computation with a wrong rounding or filter.

\paragraph{These are not retrieval failures.} \texttt{lookup\_domain} is called
on 331 of 401 tasks in \armC{}; the 70 that skip it are \emph{easier for every
arm} (\armS{} scores 41\% there against 19\% elsewhere), and 26 of the 33
losses occur on tasks where the lookup \emph{was} performed. The agent fetches
the rule and misapplies it.

That diagnosis points at a specific remedy, and points away from another.
Hoisting rule bodies into the resident prompt would not obviously help, since
the rule was already fetched. Making the wildcard semantics \emph{executable}
--- a tool that resolves which fee rules match a transaction, rather than prose
the model reimplements in SQL each time --- addresses the observed failure
directly. Section~\ref{sec:related} returns to this, because it is exactly the
boundary between what we measured and what pre-computed semantic layers
provide.
